% Online Supplement E-source: bid-layer benchmark + cost-recall grid detail
% Migrated 2026-06-06 from sec_app09_bid_benchmark_submission.tex and
% sec_app06_forensic_sequence_submission.tex during the JLEO compression pass.
% Heavy material only; the appendix retains a compact summary + pointer.
\section{Bid-Layer Benchmark and Cost-Recall Grid: Full Detail}
\label{supp:bid_cost}

This section deposits the full construction and diagnostics behind the
full-observability bid-moment benchmark (Appendix~E,
\ref{app:bid_benchmark_submission}) and the complete cost-recall grid behind
the sequential gatekeeping frontier (\ref{app:forensic_sequence_submission}).
``Cobidder'' is an adjudication-anchored exposure label throughout---an
always-loser firm that shares at least one tender-item with a BEC-active
direct CADE defendant---and never cartel membership; the benchmark
\emph{ranks} that exposure target for forensic priority and makes no claim to
identify cartelists.

%--------------------------------------------------------------------
\subsection{Bid-Layer Data, Units, and Support}
\label{supp:bid_data}

The bid layer is built from
\texttt{v3/data/processed/bid\_level\_with\_prices.parquet}
(approximately 245~MB), one row per priced bid with columns
\texttt{c\'odigofornecedor} (firm), \texttt{numerodaoc} (OC sequence),
\texttt{c\'odigoitem} (item), and \texttt{bid\_price}; a bid is retained
only if \texttt{bid\_price IS NOT NULL AND bid\_price > 0}. The benchmark
is produced by script~31 (\texttt{31\_imhof\_full\_pipeline.R},
\texttt{arrow}\,+\,\texttt{data.table}, seed \texttt{20260430}) and a
DuckDB twin, script~49 (\texttt{49\_imhof\_incremental\_value.R}, seed
\texttt{20260501}); a dedicated validation run
(\texttt{09\_bid\_benchmark\_validation.R}, seed \texttt{20260603})
generates the audit numbers. The bid layer links to the award layer
through the tender-item key $(\texttt{numerodaoc},\texttt{c\'odigoitem})$
and, for firm attachment, through
$(\texttt{numerodaoc},\texttt{c\'odigoitem},\texttt{c\'odigofornecedor})$;
the within-tender moments are \emph{not} recoverable from the award layer,
which records no per-bid prices.

The modeling unit is the \emph{firm}: per-tender moments are shared across
all firms bidding on a tender-item and then firm-mean aggregated. Of the
$16{,}843$ always-losers, $64$ are dropped because they have no priced
bids or fewer than two bids per tender (so the coefficient of variation,
skewness, and kurtosis are undefined). Excluding the direct CADE
defendants by upstream label curation leaves an evaluation pool of
$16{,}731$ firms, of which $\valBidPosRetained$ are always-loser cobidder
positives. All $\valBidPosRetained$ positive (cobidder) labels are
non-missing on every Imhof moment: the support loss falls entirely on the
negatives, so there is no positive attrition.

%--------------------------------------------------------------------
\subsection{Bid-Feature Dictionary}
\label{supp:bid_dictionary}

Seven Imhof--Wallimann moments are constructed within each tender-item
with $\geq 2$ priced bids, then averaged across each firm's priced tenders
to the firm level. They fall into four families: \emph{dispersion}
(coefficient of variation \texttt{imhof\_cv\_mean} and its within-firm SD
\texttt{imhof\_cv\_sd}, undefined for single-tender firms);
\emph{level/relative spread} (normalized spread and log min--max,
\texttt{imhof\_spread\_mean}, \texttt{imhof\_minmax\_mean});
\emph{shape} (skewness \texttt{imhof\_skew\_mean} and excess kurtosis
\texttt{imhof\_kurt\_mean}); and \emph{rank/order} (the second-low gap
$\log(b_2/b_1)$, \texttt{imhof\_second\_low\_mean}, the only moment with a
winning-margin component, of low--moderate leakage risk but still a
field-level moment shared across firms, not a firm-specific outcome).
Two auxiliary counts accompany them, \texttt{n\_bids} (gating the $\geq 2$
filter) and \texttt{n\_tenders\_priced} (a firm exposure count). No
co-bidding/network or timing, bid-revision, or late-bid features are
implemented: the parquet read selects only firm, OC, item, and price, and
no timestamp column is touched.

%--------------------------------------------------------------------
\subsection{Model Specification and Validation Designs}
\label{supp:bid_model}

The learner is a random forest (\texttt{ranger}, $500$ trees,
\texttt{probability = TRUE}, $\texttt{num.threads}=12$)---not a logit,
despite an earlier script header that mislabeled it. There is no
hyperparameter tuning (fixed $500$ trees, default \texttt{mtry} and
\texttt{min.node.size}, no grid search) and no standardization (the forest
is scale-invariant). Missingness is handled listwise through the pool
restriction. The outcome is binary \texttt{is\_cade} ($1$ for an
always-loser cobidder in the candidate pool, not a direct defendant).
Evaluation is by $5$-fold cross-validation, each firm receiving one
held-out predicted probability; AUC and DeLong tests are computed on these
out-of-fold firm-level predictions.

Four validation designs are run, increasing in stringency:
\textbf{Design~A} (pooled-random $5$-fold) is a diagnostic
\emph{upper bound}; \textbf{Design~E} (case-grouped / leave-one-case-out)
holds out an entire CADE case so its clustered positives never co-train;
\textbf{Design~F} (exclude-label-defining tenders) recomputes the moments
after dropping defendant-co-appearance tenders to test feature--target
contamination; and \textbf{Design~G} (environment holdout) holds out an
item-group. Strict bid-timing validation is \textbf{blocked}: the
within-tender moments cannot be cleanly time-limited without re-deriving
them on a per-period bid panel. Because positives cluster by case, random
folds (Design~A) are optimistic---same-case positives appear in both
train and test.

%--------------------------------------------------------------------
\subsection{Bid-Layer Performance by Validation Design}
\label{supp:bid_perf}

The award score is \emph{clean} (a CADE-label-independent rank on
\texttt{log1p(tenders\_count)}), as is its strict-timing variant. The bid
and combined forests under full-period features and random folds are
\textbf{diagnostic-only}: random folds are optimistic, and the firm-mean
moments include label-defining tenders. The contamination is small:
excluding label-defining tenders (Design~F) leaves the bid ROC essentially
unchanged ($0.717 \to 0.713$), so the fragility is case clustering, not
feature--target overlap. The case-grouped result (Design~E) and the
contamination check (Design~F) together establish that the bid layer's
apparent reach is driven by case clustering rather than by feature--target
overlap.

\begin{table}[htbp]\centering
\begin{adjustbox}{max width=\textwidth,center}
\begin{threeparttable}
\caption{Bid-layer performance by validation design.}
\label{supp:tab_bid_perf}\scriptsize
\begin{tabular}{llrrr}
\toprule
Design & Model & ROC-AUC & PR-AUC & Prec@500 \\
\midrule
A (random)        & Award (cont.) & 0.760 & 0.143 & 0.216 \\
A (random)        & Award FL14    & 0.688 & 0.098 & 0.130 \\
A (random)        & Bid RF        & 0.717 & 0.116 & 0.180 \\
A (random)        & Combined RF   & 0.756 & 0.188 & 0.274 \\
E (case-grouped)  & Award (cont.) & 0.760 & 0.143 & 0.216 \\
E (case-grouped)  & Award FL14    & 0.688 & 0.098 & 0.130 \\
E (case-grouped)  & Bid RF        & 0.626 & 0.062 & 0.086 \\
E (case-grouped)  & Combined RF   & 0.689 & 0.103 & 0.190 \\
F (excl.\ label)  & Award (cont.) & 0.829 & 0.156 & 0.216 \\
F (excl.\ label)  & Bid RF        & 0.713 & 0.088 & 0.130 \\
F (excl.\ label)  & Combined RF   & 0.855 & 0.217 & 0.302 \\
G (env.\ holdout) & Award (cont.) & 0.768 & 0.185 & 0.150 \\
G (env.\ holdout) & Bid RF        & 0.747 & 0.171 & 0.150 \\
G (env.\ holdout) & Combined RF   & 0.734 & 0.212 & 0.144 \\
\bottomrule
\end{tabular}
\begin{tablenotes}\scriptsize
\item $N=16{,}731$ always-losers with $\valBidPosRetained$ always-loser
cobidder positives (defendant-excluded pool) for Designs~A, E; $N=16{,}665$
with $569$ positives for Design~F; $N=2{,}377$ with $129$ positives for the
held-out environment in Design~G. PR-AUC is reported alongside ROC because
the positive base rate is near one percent. Design~A is a
full-observability \emph{diagnostic upper bound} with random-CV optimism;
the case-grouped block (E) is the honest robustness test, where the
combined model's PR-AUC ($0.103$) falls below award-only ($0.143$). The
award rank is identical across A and E by construction. Cobidders are
adjudication-anchored exposure, not cartel membership; false positives are
unlabeled firms, not ``bad firms.'' Lift omitted; PR-AUC is the
rare-target metric and carries the comparison.
\end{tablenotes}
\end{threeparttable}
\end{adjustbox}
\end{table}

%--------------------------------------------------------------------
\subsection{Full Cost-Recall Grid and Frontier Figure}
\label{supp:cost_grid}

Table~\ref{supp:tab_G_frontier} reports a slice of the cost-recall grid
for the sequential award\,$\to$\,bid rule at the $k=500$ forensic-queue
depth, across the first-stage cut $K_1$; the full machine-readable grid
spans all rules, all $K_1$, and $k \in \{100,250,500,1000\}$ with every
cost denominator.

\begin{table}[htbp]\centering
\begin{adjustbox}{max width=\textwidth,center}
\begin{threeparttable}
\caption{Cost-recall slice, sequential award\,$\to$\,bid, queue depth $k=500$.}
\label{supp:tab_G_frontier}
\scriptsize
\begin{tabular}{rrrrrr}
\toprule
$K_1$ & Firms opened & Bid rows opened & Bid-row red.\ & TP & Recall \\
\midrule
$500$    & $500$    & $1{,}356{,}346$ & $0.60$ & $108$ & $0.166$ \\
$1{,}000$ & $1{,}000$ & $1{,}783{,}750$ & $0.47$ & $124$ & $0.190$ \\
$2{,}000$ & $2{,}000$ & $2{,}283{,}924$ & $0.33$ & $116$ & $0.178$ \\
$3{,}000$ & $3{,}000$ & $2{,}575{,}274$ & $0.24$ & $115$ & $0.177$ \\
$5{,}000$ & $5{,}000$ & $2{,}921{,}977$ & $0.14$ & $102$ & $0.157$ \\
$16{,}772$ & $16{,}772$ & $3{,}393{,}915$ & $0.00$ & $102$ & $0.157$ \\
\bottomrule
\end{tabular}
\begin{tablenotes}\scriptsize
\item \textit{Source:} \texttt{table\_6\_cost\_recall\_frontier.csv}
(sequential award$\to$bid, $k=500$). ``Bid-row red.'' is the
fraction of the full-pool bid-row footprint avoided,
$1-\text{(bid rows opened)}/3{,}393{,}915$. TP counts
adjudication-anchored positives in the top-$k$ queue (pool $=\valCostNpos$).
\item TP is non-monotone in $K_1$ because the bid rerank reorders
survivors within a fixed queue depth: $K_1=1{,}000$ recovers more true
positives ($124$) than $K_1=2{,}000$ ($116$) at $k=500$, one more sign
that no $K_1$ is optimal. The joint full-observability TP at $k=500$ is
$133$, which the sequential rule approaches ($124/133 = 93\%$ at
$K_1=1{,}000$) but does not exceed.
\end{tablenotes}
\end{threeparttable}
\end{adjustbox}
\end{table}

Figure~\ref{supp:fig_G_cost_recall_ti} plots recall against the tender-item
footprint for the award-only, sequential, and joint rules, making the
frontier shape visible: the sequential rule tracks close to the joint
upper bound while opening fewer tender-items, but never crosses it.

\begin{figure}[htbp]\centering
\includegraphics[width=0.86\textwidth]{./output/figures/fig_G_cost_recall_tender_items}
\caption{Recall versus opened tender-items, by ranking rule.}
\label{supp:fig_G_cost_recall_ti}
\par\medskip
{\footnotesize\textit{Alt text:} A cost-recall scatter with opened
tender-items on the horizontal axis and adjudication-anchored recall on
the vertical axis. The joint full-observability rule sits at the far
right (largest footprint) at the top of the recall range; award-only
points sit at the far left (smallest footprint) at lower recall;
sequential award\,$\to$\,bid points trace an intermediate frontier that
approaches but does not exceed the joint recall while opening fewer
tender-items.}
\end{figure}

%--------------------------------------------------------------------
\subsection{Baselines and Case-Holdout Concentration}
\label{supp:cost_holdout}

The sequential rule sits well above a same-size random survivor floor
across the grid, confirming that it is the award ranking, not the mere act
of opening $K_1$ firms, that carries the signal. The award-ranked survivor
pool captures most of the $\valCostNpos$ positives \emph{before} any bid
rerank, so the bid rerank reorders that pool into a forensic-priority queue
rather than recovering positives the award gate missed. Because adjudicated
positives cluster in a handful of CADE cases, apparent recall can be
carried by a single large case: the largest case supplies $32.0\%$ of the
positives and $45.4\%$ of the true positives in the top-$500$ queue.
Grouping firms by case and holding out the largest case lowers recall on
the remaining smaller cases materially, with several individual small cases
near zero. Performance is therefore case-fragile; the transferable object
is the decomposition framework, not a portable cartel score.

A timing-disciplined gate must rank on information available before the
investigation window. Re-scoring the award gate on a strict pre-period
(2009--2016) and evaluating on later positives, its strict-timing
discrimination is AUC $0.684$ on the training-AL pool against $0.761$ for
the full-window score---the honest cost of respecting the timing order.
The sequential rule cannot be subjected to the same test, because the
within-tender moments cannot be re-derived on a pre-period basis without
re-pipelining the raw LANCES bid records, so we report it as infeasible
rather than estimate it on contaminated features.

%--------------------------------------------------------------------
\subsection{Gatekeeping Algorithm and Operating-Point Logic}
\label{app:algorithm_submission}

No single $(K_1,k)$ pair is preferred for every agency; the right point
depends on capacity and on whether the priority is recall or precision,
so the frontier is a menu, not a cutoff. The operating-point grid with
recall, precision, and recovered-positive counts at each $K_1$ is the
main text's Table~6. The full gatekeeping procedure can be stated
compactly.

\begin{quote}
\textbf{Algorithm E.1: Award-Layer Gatekeeping with Cost Measurement}
\begin{enumerate}
\item Define the firm universe from award records (participant identity,
winner identity, item code, negotiated price, procuring unit, year,
modality).
\item For each firm $i$, compute the win count $W_i$ and the
participation count $T_i$ over the calibration window; restrict to the
zero-win stratum $\{i : W_i = 0\}$.
\item Compute the award score $s_i = \log(1+T_i)$.
\item Over a grid of first-stage cuts $K_1$, form the survivor pool
$\mathcal{S}_{K_1}$ of the top-$K_1$ firms (or apply the
median-plus-$1.5\times$IQR threshold, or a capacity budget).
\item Recover bid microdata only for the tender-items touched by
$\mathcal{S}_{K_1}$, and record the cost denominators D1--D8 at that
$K_1$.
\item Compute the bid-layer forensic features inside $\mathcal{S}_{K_1}$.
\item Rerank survivors by the bid score alone or by the award-plus-bid
combined score, over a grid of queue depths $k$.
\item Deliver the top-$k$ queue from that ranking as a forensic-priority
list for investigative review.
\item Trace the cost-recall frontier over $(K_1,k)$ and select the
operating point that fits the agency's capacity and recall/precision
priority; treat the output as forensic priority, not as a liability
finding.
\end{enumerate}
\end{quote}

Joint scoring requires bid microdata for the full pool before triage;
sequential gatekeeping requires it only after the award-layer survivor
pool is formed. The survivor-pool size, the threshold or top-$K_1$ choice,
the queue depth $k$, the bid-versus-combined rerank, and the
recalibration cadence are institutional choices the algorithm leaves to
the agency.
